\section{The fall of a good creature}
\label{sec:angel-fall}

\subsection{A good desired against its proper rule (Question 63)}

How could a lucid intellect choose evil without bodily passion or an already corrupt habit? Aquinas's first answer is to distinguish the goodness of the object from the rightness of the act. Something desirable in itself can be sought without the measure and order it requires. Created excellence is good; possessing it as though one's own will were its ultimate rule is not. A creature can turn from the divine rule without first mistaking evil for a good nature. The fault can consist in failing to attend to what should regulate the choice.\footnote{\ST{I q.~63 a.~1}, especially ad~4; \emph{De malo} q.~16 a.~2, corpus.}

Augustine makes the distinction tangible by considering gold and power. Gold does not become an evil substance because avarice prefers it to justice; authority does not become an evil substance because pride refuses a higher rule. The disorder belongs to the will's preference. Asking for a positive substance that manufactured the first evil will therefore mistakes a failure for a new created nature. Augustine compares it to trying to see darkness as a positive color. The defect remains voluntary: a spirit departs from the greater good through its own willing. This joins metaphysics to moral responsibility without giving evil a creator alongside God.\footnote{Augustine, \emph{City of God} XII.7--8.}

Aquinas consequently identifies pride as the first angelic sin. Envy follows: another's excellence appears an obstacle to the singular excellence the proud creature wants. Augustine had already refused to blame bodily flesh for every sin, since the devil's pride and envy do not arise from sensual indulgence. Demons can be guilty of inciting bodily sins in others without themselves possessing the bodily appetites proper to those sins.\footnote{\ST{I q.~63 a.~2}; Augustine, \emph{City of God} XIV.3.}

Article 3 asks what it means to desire to be like God. The desire cannot mean an angel's reasonable expectation of becoming the infinite Creator. Its natural knowledge excludes that impossibility. Nor is every desire for likeness to God sinful: the creature is made for such likeness. Aquinas instead identifies a refusal of the right mode of receiving perfection. He gives two closely related accounts: choosing natural self-sufficiency as final happiness instead of supernatural beatitude, or seeking supernatural beatitude as attainable by one's own power instead of God's gracious gift. The \emph{De malo} concentrates the explanation on the second: the devil would receive God's causality in nature but refuses dependence on grace for the final end.\footnote{\ST{I q.~63 a.~3}; \emph{De malo} q.~16 a.~3, corpus.}

Article 4 excludes an evil angelic substance. God made these creatures good; their wickedness belongs to voluntary defection. Lateran IV teaches precisely that distinction. The devil is neither a second first principle nor a nature God created evil.\footnote{\ST{I q.~63 a.~4}; Lateran IV, \emph{Firmiter}, DS~800; \emph{CCC} 391--392.}

The next two articles concern the first instants, not a chronology recoverable from a human calendar. Aquinas denies sin in the very first instant of creation. He does not claim that free choice needs a long deliberation: an angel can choose immediately. His argument in the \emph{Summa} concerns the good Creator's initial causality. In the \emph{De malo}, he develops instead the order between connatural self-knowledge and movement towards a supernatural end. These are different lines of explanation for the same conclusion. He then regards the fall immediately after the first instant as more probable, given his positions on creation in grace and immediate reward after merit.\footnote{\ST{I q.~63 aa.~5--6}; \emph{De malo} q.~16 a.~4, corpus. Aquinas himself records alternative accounts.}

The hardest detail is article 5's fourth reply: on the hypothesis of creation in grace, all initially merited, but some immediately placed an obstacle to the beatitude they had merited. This does not mean a blessed angel fell out of the vision of God. The distinction is between an act under grace and the consummated gift that makes subsequent sin impossible. Nor does Aquinas offer these ``instants'' as pieces of bodily time: article 6 explicitly reckons the succession of intellectual and voluntary acts.\footnote{\ST{I q.~63 a.~5 ad~4; a.~6 ad~4}; compare \ST{I q.~62 a.~8}.}

The final three articles address leadership and multitude. Aquinas favors Gregory's view that the chief rebel was the highest angel, since excellence can furnish the occasion for pride. He nevertheless permits Damascene's account of a ruler of the earthly order. The leader causes others' sins by exhortation and consent, not by compelling their wills. Finally Aquinas holds that more angels remained faithful than fell. The dragon drawing a third of the stars supplies article 8 with its image of a leader drawing others into rebellion. Article 9 concludes with the greater multitude of the faithful, without fixing an exact proportion.\footnote{\ST{I q.~63 aa.~7--9}; Damascene, \emph{Orthodox Faith} II.4; Apoc.~12:4; 4~Kings~6:16.}

\angeldossier{I q.~63, aa.~1--9}{Defined: the devil and demons were created good and became evil by themselves (Lateran IV, DS~800). Common teaching: pride and voluntary rebellion. Thomist position or disputed account: the precise object, first instants, original rank, and relative numbers.}{Augustine, \emph{City of God} XIV.3; Damascene II.4; Isaiah~14:12--14 and Ezechiel~28 in Aquinas's angelic reading; Gregory, Gospel homily~34, invoked in a.~7; \emph{De malo} q.~16 aa.~2--4.}{Damascene's earthly ruler differs from the highest-angel account preferred by Aquinas. The alternatives concerning grace and first acts condition Aquinas's proposed sequence.}

\subsection{Knowledge without wisdom, regret without repentance (Question 64)}

Punishment does not make a demon cease to be an intellectual creature. Article 1 distinguishes natural knowledge, revealed information, and the loving wisdom that unites a knower to God. Natural intellectual endowment remains. Some divine matters can still be disclosed as God's government requires. The loving knowledge of wisdom is absent. Thus genuine information and spiritual ruin can coexist: possessing truths is not the same as possessing the good those truths disclose.\footnote{\ST{I q.~64 a.~1}, corpus and ad~4--5.}

This explains both the extent and the limits of demonic knowledge. Aquinas denies that the demons ever possessed the beatific vision. Their knowledge of Christ came through temporal effects and was incomplete. When a previously uncertain event occurs, they can know it as present: their new experience does not mean the development of sensory organs or omniscience. Their intelligence remains finite, and their opposition to divine wisdom can corrupt judgments about supernatural action.\footnote{\ST{I q.~64 a.~1 ad~4--5}; \ST{I q.~58 a.~5}.}

Article 2 asks why the choice cannot be revoked. Aquinas appeals to the settled manner of angelic apprehension, unlike our successive reconsideration. The \emph{De malo} supplies a helpful qualification: freedom to choose different actions is not identical with freedom to revise the fundamental orientation already fixed. Demons can choose diverse means while all remain governed by their first aversion. Their state as wayfarers has ended. John Damascene compares angelic fall to human death, and the Catechism locates the irreversibility in the character of their choice, not a deficiency of divine mercy.\footnote{\ST{I q.~64 a.~2}; \emph{De malo} q.~16 a.~5, corpus; Damascene, \emph{Orthodox Faith} II.4; \emph{CCC} 393.}

Sorrow therefore need not mean repentance. Article 3 distinguishes willing against a punishment from repudiating one's fault because it offends God. The demons lack bodily passions but can will that an unwelcome reality were otherwise. Their frustration bears witness to the goods their nature still seeks, especially happiness; it does not by itself heal their chosen hostility.\footnote{\ST{I q.~64 a.~3}, especially ad~3.}

Article 4 locates their punishment within providence and final judgment. Aquinas distinguishes hell as their due punishment from the present sphere in which their permitted opposition can test human beings. His ``dark air'' belongs to the biblical and medieval cosmological description; the theological claim is that their activity remains bounded by divine government. A spirit suffers confinement through an unwilling subjection, not through bodily compression. Their present activity does not suspend their punishment, just as a good angel's mission does not interrupt its blessedness.\footnote{\ST{I q.~64 a.~4}, especially ad~1 and ad~3.}

\angeldossier{I q.~64, aa.~1--4}{Church teaching: the irrevocable angelic fall (\emph{CCC} 393). Common teaching: continuing punishment and loss of blessedness. Thomist position: the psychological account of fixation and sorrow; disputed cosmological account: the ``dark air.''}{Dionysius, \emph{Divine Names} IV; Damascene II.4; Matthew~25:41, 46; James~2:19; Augustine, \emph{City of God} IX.21 and XIX.13, as invoked by Aquinas; \emph{De malo} q.~16 a.~5.}{Aquinas opposes the prospect of recurrent falls and restorations he attributes to Origen. He also rejects postponing every punishment until the final judgment.}
